%% ========================================================================
%% Chapter 11: File & User/Group Management
%% ========================================================================

\chapter{File \& User/Group Management}
\label{ch:file-and-user-management}

\begingroup
\small
\setstretch{0.85}
\setlength{\parskip}{0pt}
\setlist{nosep}

This chapter summarizes the foundations of access administration in Linux: classic Unix
permissions, ACLs for more granular cases, user and group management, using \cmd{sudo}
and \cmd{su}, and disk space restrictions with quota. The focus of this chapter is to
understand Linux's \emph{access control model} as a whole: who owns files, who can
access them, how administrative rights are delegated, and how resource usage is limited
in multi-user systems.

\textbf{Learning Objectives}

After completing this chapter, readers are expected to be able to:

\begin{enumerate}
    \item read and modify file/directory permissions;
    \item distinguish the roles of owner, group, and others;
    \item use ACLs when standard permissions are insufficient;
    \item create, modify, and delete users and groups;
    \item manage administrative rights with \cmd{sudo} securely; and
    \item enable and check disk quota for users or groups.
\end{enumerate}

%% ========================================================================
\section{Access Control System (Permissions)}
\label{sec:permissions}
%% ========================================================================

\textbf{Concept Explanation}

Linux inherits the simple yet powerful Unix access control model. Each file and
directory always has three access identities:

\begin{description}
    \item[Owner] the user who owns the file.
    \item[Group] the group associated with the file.
    \item[Others] all other users outside the two entities above.
\end{description}

For each of these entities, three basic permissions are available: read, write, and execute.

\begin{table}[htbp]
    \centering
    \caption{Meaning of basic Linux permissions}
    \label{tab:permissions-type}
    \begin{tblr}{
        colspec = {lX[l]X[l]},
        width = \linewidth,
        hlines,
        vlines,
        row{1} = {font=\bfseries},
        cells = {font=\small}
    }
    \textbf{Symbol} & \textbf{On files} & \textbf{On directories} \\
    \texttt{r} & Read file contents & List directory contents \\
    \texttt{w} & Modify file contents & Create or delete entries in directory \\
    \texttt{x} & Execute file & Enter directory \\
    \end{tblr}
\end{table}

\begin{notebox}
The \texttt{w} right on a file only modifies file contents. Deleting a file is determined
by permissions on the directory containing the file.
\end{notebox}

\textbf{Commands Used}

\begin{itemize}
    \item \texttt{ls -l}: view file type, permissions, owner, and group.
    \item \texttt{chmod}: change file or directory permissions.
    \item \texttt{chown}: change owner and/or group of a file.
    \item \texttt{chgrp}: change group owner of a file.
    \item \texttt{umask}: view or change default permissions for new files.
\end{itemize}

\subsection{Reading \texttt{ls -l} Output}

The \cmd{ls -l} output summarizes file type, permissions, owner, group, size, and time.

\begin{lstlisting}[language=bash, caption={Reading permissions with ls -l}]
$ ls -l /etc/passwd
-rw-r--r-- 1 root root 2847 Jan 15 10:30 /etc/passwd
\end{lstlisting}

Permission string interpretation:

\begin{itemize}
    \item first character: file type (\texttt{-}, \texttt{d}, \texttt{l}, etc.)
    \item next three characters: owner permissions
    \item next three: group permissions
    \item last three: others permissions
\end{itemize}

Example \texttt{-rwxr-x---} means regular file, owner has full rights, group can
read and execute, others have no access.

\subsection{Symbolic and Octal Notation}

Permissions can be written in two forms:

\begin{itemize}
    \item \textbf{symbolic}: relative, e.g., \texttt{u+x}, \texttt{go-w}
    \item \textbf{octal}: absolute, e.g., \texttt{755}, \texttt{640}, \texttt{600}
\end{itemize}

\begin{table}[htbp]
    \centering
    \caption{Permission conversion to octal notation}
    \label{tab:octal-permissions}
    \begin{tblr}{
        colspec = {lccX[l]},
        width = \linewidth,
        hlines,
        vlines,
        row{1} = {font=\bfseries},
        cells = {font=\small}
    }
    \textbf{Permission} & \textbf{Binary} & \textbf{Octal} & \textbf{Meaning} \\
    \texttt{---} & 000 & 0 & No permissions \\
    \texttt{r--} & 100 & 4 & Read only \\
    \texttt{rw-} & 110 & 6 & Read and write \\
    \texttt{rwx} & 111 & 7 & Read, write, execute \\
    \end{tblr}
\end{table}

Most common patterns:

\begin{itemize}
    \item \texttt{755}: public directories or programs
    \item \texttt{644}: regular files
    \item \texttt{600}: private files like SSH keys
    \item \texttt{700}: private scripts or directories
\end{itemize}

\subsection{\cmd{chmod}, \cmd{chown}, and \cmd{chgrp}}

Three core commands for permission management are:

\begin{itemize}
    \item \cmd{chmod}: change permissions
    \item \cmd{chown}: change owner and/or group
    \item \cmd{chgrp}: change group only
\end{itemize}

\begin{lstlisting}[language=bash, caption={Basic permission and ownership commands}]
# symbolic chmod
chmod +x script.sh
chmod u=rw,g=r,o= file.txt
chmod go-w file.txt

# octal chmod
chmod 644 /etc/app.conf
chmod 755 /usr/local/bin/my-script.sh
chmod 600 ~/.ssh/id_rsa

# ownership
sudo chown alice report.txt
sudo chown alice:developers report.txt
sudo chgrp developers report.txt
\end{lstlisting}

Usage principles:

\begin{itemize}
    \item use symbolic notation for small, incremental changes;
    \item use octal notation when you want definite results, especially in scripts;
    \item always verify recursive changes with \cmd{ls -la} on several files.
\end{itemize}

\begin{warningbox}
Avoid \texttt{chmod 777} unless you fully understand the risks. This permission
gives write and execute rights to all users.
\end{warningbox}

\subsection{Special Permission Bits}

Besides the nine \texttt{rwx} bits, Linux has three special bits:

\begin{description}
    \item[SUID] program runs with the UID of the file owner. Classic example:
        \cmd{/usr/bin/passwd} needs to write to \file{/etc/shadow}.
    \item[SGID] on executable files, the process runs with the file's GID; on directories,
        new files inherit the directory's group.
    \item[Sticky bit] on directories, users can only delete their own files even if
        the directory is writable by many people.
\end{description}

\begin{lstlisting}[language=bash, caption={Special bits examples}]
# SUID
ls -l /usr/bin/passwd
sudo chmod 4755 /usr/local/bin/my-tool

# SGID on shared directory
sudo mkdir /srv/project
sudo chown root:developers /srv/project
sudo chmod 2775 /srv/project

# Sticky bit
ls -ld /tmp
sudo chmod 1777 /shared/uploads
\end{lstlisting}

\begin{tipbox}
The combination of SGID + sticky bit is often used for collaborative directories:
SGID keeps the group consistent, sticky bit prevents mutual file deletion.
\end{tipbox}

\subsection{umask --- Default New Permission Value}

\textbf{umask} is a mask that \emph{reduces} default permissions for new files.
Practically, regular files start from \texttt{666}, directories from \texttt{777},
then reduced by the umask value.

\begin{table}[htbp]
    \centering
    \caption{umask effect on new files and directories}
    \label{tab:umask-effect}
    \begin{tblr}{
        colspec = {lcccc},
        width = \linewidth,
        hlines,
        vlines,
        row{1,2} = {font=\bfseries},
        cells = {font=\small}
    }
    \SetCell[r=2]{c} \textbf{umask} & \SetCell[c=2]{c} \textbf{New file} & \SetCell[c=2]{c} \textbf{New directory} \\
    & Octal & Symbolic & Octal & Symbolic \\
    \texttt{022} & \texttt{644} & \texttt{rw-r--r--} & \texttt{755} & \texttt{rwxr-xr-x} \\
    \texttt{027} & \texttt{640} & \texttt{rw-r-----} & \texttt{750} & \texttt{rwxr-x---} \\
    \texttt{077} & \texttt{600} & \texttt{rw-------} & \texttt{700} & \texttt{rwx------} \\
    \end{tblr}
\end{table}

\begin{lstlisting}[language=bash, caption={Viewing and changing umask}]
umask
umask -S
umask 027
echo "umask 027" >> ~/.bashrc
\end{lstlisting}

\begin{notebox}
umask never adds execute to regular files. If a file must be executable, use
\cmd{chmod} explicitly.
\end{notebox}

\subsection*{Lab 9.1 --- Permissions}

\textbf{Objective}: read permissions, change them safely, and understand the effects
of \cmd{chmod}, \cmd{chown}, and \cmd{umask} on real files.

\textbf{Step 1}: Create working directory and two test files.

\begin{lstlisting}[language=bash, caption={Preparing permission test files}]
mkdir ~/lab-permissions && cd ~/lab-permissions
echo "secret data" > secret.txt
echo '#!/bin/bash' > myscript.sh
echo 'echo Hello' >> myscript.sh
ls -la
\end{lstlisting}

Note the initial permissions of both files. Typically new files don't have the execute bit.

\textbf{Step 2}: Make \file{secret.txt} private to owner only.

\begin{lstlisting}[language=bash, caption={Making file private}]
chmod 600 secret.txt
ls -l secret.txt
\end{lstlisting}

Permission \texttt{600} means owner can read and write, while group and others have no access.

\textbf{Step 3}: Make \file{myscript.sh} executable.

\begin{lstlisting}[language=bash, caption={Making script executable}]
chmod 755 myscript.sh
ls -l myscript.sh
./myscript.sh
\end{lstlisting}

Permission \texttt{755} gives full rights to owner, and read+execute rights to group
and others. Without the execute bit, script files cannot be run directly.

\textbf{Step 4}: Create shared directory and observe simple SGID effect.

\begin{lstlisting}[language=bash, caption={Shared directory with SGID}]
mkdir shared-dir
chmod g+s shared-dir
ls -ld shared-dir
\end{lstlisting}

If output shows letter \texttt{s} in the group execute position, SGID is active.

\textbf{Step 5}: Test \cmd{umask} effect on new files.

\begin{lstlisting}[language=bash, caption={Observing umask effect}]
umask
umask 027
touch testfile-027
ls -l testfile-027
\end{lstlisting}

\textbf{Analysis}

\begin{enumerate}
    \item Why can't \file{secret.txt} be read by group and others after \cmd{chmod 600}?
    \item What's the difference in meaning between \texttt{600} and \texttt{755} for the tested files?
    \item After \cmd{umask 027}, what permissions are produced for new files, and why not \texttt{777}?
\end{enumerate}

\textbf{Challenge}

Change the owner or group of one test file to another account or group available on
the system, then explain the changes in \cmd{ls -l} output before and after.

%% ========================================================================
\section{Access Control Lists (ACLs)}
\label{sec:acls}
%% ========================================================================

Classic Unix permissions work well for most cases, but they only recognize one owner
and one group. Problems arise when access needs to be given to one specific user
without changing the file owner and without adding that user to the main group.

Common problem example:

\begin{examplebox}[Why ACL is needed]
File \file{report.txt} is owned by \texttt{alice} and group \texttt{finance}. User
\texttt{bob} needs to read the file, but you don't want all other group members to
get access. ACL allows you to give permission only to \texttt{bob}.
\end{examplebox}

ACL adds \emph{additional} access rules on top of standard permissions. There are two types:

\begin{description}
    \item[Access ACL] applies directly to the current file or directory.
    \item[Default ACL] for directories only, inherited by new content inside.
\end{description}

\begin{notebox}
In Ubuntu 22.04, modern filesystems like ext4, XFS, and Btrfs generally support ACL.
If utilities are not available, install the \texttt{acl} package.
\end{notebox}

\textbf{Commands Used}

\begin{itemize}
    \item \texttt{getfacl file}: display complete ACL of file or directory.
    \item \texttt{setfacl -m ...}: add or modify ACL entry.
    \item \texttt{setfacl -x ...}: remove one ACL entry.
    \item \texttt{setfacl -b file}: remove all additional ACLs.
    \item \texttt{setfacl -d -m ... dir}: set \textit{default ACL} on directory.
\end{itemize}

\subsection{Step 1: Reading Existing ACLs}

The first command to know is \cmd{getfacl}. Its function is similar to \cmd{ls -l},
but more detailed.

\begin{lstlisting}[language=bash, caption={Reading basic ACL with getfacl}]
sudo apt install -y acl
getfacl file.txt

# minimal output example
# user::rw-
# group::r--
# other::r--
\end{lstlisting}

If a file doesn't have additional ACLs yet, \cmd{getfacl} output only shows three
basic entries equivalent to regular Unix permissions: owner, group, and others.
Once additional ACLs are applied, you'll see specific user/group entries and often
a \texttt{mask}.

The most important entry formats:

\begin{table}[htbp]
    \centering
    \caption{Common ACL entry formats}
    \label{tab:setfacl-entries}
    \begin{tblr}{
        colspec = {lX[l]X[l]},
        width = \linewidth,
        hlines,
        vlines,
        row{1} = {font=\bfseries},
        cells = {font=\small}
    }
    \textbf{Entry} & \textbf{Example} & \textbf{Meaning} \\
    \texttt{u::perm} & \texttt{u::rw-} & Owner permission \\
    \texttt{u:name:perm} & \texttt{u:bob:r--} & ACL for specific user \\
    \texttt{g::perm} & \texttt{g::r--} & Standard group permission \\
    \texttt{g:name:perm} & \texttt{g:audit:rw-} & ACL for specific group \\
    \texttt{o::perm} & \texttt{o::---} & Others permission \\
    \texttt{m::perm} & \texttt{m::rw-} & ACL mask \\
    \end{tblr}
\end{table}

When ACL is active, \cmd{ls -l} displays a \texttt{+} sign at the end of the
permission string, e.g., \texttt{-rw-r-----+}.

\subsection{Step 2: Adding ACL for One User or Group}

The main command for creating or modifying ACL is \cmd{setfacl}. Start from the
simplest case: giving read access to one user.

\begin{lstlisting}[language=bash, caption={Adding ACL gradually}]
# give bob read permission
setfacl -m u:bob:r file.txt

# give audit group read-write permission
setfacl -m g:audit:rw file.txt

# view results
getfacl file.txt
\end{lstlisting}

How to read it:

\begin{itemize}
    \item \texttt{-m} means \emph{modify}: add or change entry;
    \item \texttt{u:bob:r} means user \texttt{bob} gets read permission;
    \item \texttt{g:audit:rw} means group \texttt{audit} gets read and write permission.
\end{itemize}

If ACL is no longer needed, remove one entry or all additional entries:

\begin{lstlisting}[language=bash, caption={Removing ACL}]
setfacl -x u:bob file.txt
setfacl -x g:audit file.txt
setfacl -b file.txt
\end{lstlisting}

\texttt{-x} removes one specific entry. Option \texttt{-b} removes all additional
ACLs and returns the file to standard Unix permission model.

\subsection{Step 3: Understanding Mask and Effective Permissions}

Once additional ACL is installed, Linux usually creates a \textbf{mask}. For beginners,
think of mask as the \emph{upper limit} of permissions that can actually apply to
certain ACL entries.

\begin{lstlisting}[language=bash, caption={ACL mask and effective permissions}]
getfacl ~/report.txt
# user::rw-
# user:charlie:rwx        #effective:rw-
# group::r--              #effective:r--
# group:audit:rw-         #effective:rw-
# mask::rw-
# other::---

setfacl -m m::r-- ~/report.txt
chmod g-w ~/report.txt
\end{lstlisting}

In the example above, \texttt{user:charlie:rwx} doesn't actually get \texttt{rwx}
because the mask only allows \texttt{rw-}. That's why \cmd{getfacl} displays the
comment \texttt{\#effective:rw-}.

Practical rules to remember:

\begin{itemize}
    \item owner and others are not filtered by mask;
    \item named user, standard group, and named group can be restricted by mask;
    \item \cmd{chmod} on ACL files often affects the mask;
    \item if results are confusing, always check again with \cmd{getfacl}.
\end{itemize}

\subsection{Step 4: Default ACL for Directories}

After understanding ACL on one file, the next step is automatically inherited ACL.
This is the function of \textbf{default ACL}. It's very useful for shared project
directories.

\begin{lstlisting}[language=bash, caption={Default ACL on collaborative directory}]
mkdir ~/project
setfacl -d -m g:developers:rwx ~/project
setfacl -d -m u:deploy:r-x,g:auditors:r-- ~/project

touch ~/project/newfile.txt
getfacl ~/project/newfile.txt

mkdir ~/project/subdir
getfacl ~/project/subdir

setfacl -x d:g:developers ~/project
setfacl -k ~/project
\end{lstlisting}

Workflow step by step:

\begin{enumerate}
    \item add default ACL to directory with option \texttt{-d -m};
    \item create new file or subdirectory;
    \item check inherited ACL result with \cmd{getfacl};
    \item remove default ACL when no longer needed.
\end{enumerate}

\begin{notebox}
New files don't inherit execute from default ACL if they're regular files;
subdirectories still retain execute because \texttt{x} on directories means entry permission.
\end{notebox}

\subsection*{Lab 9.2 --- ACL}

\textbf{Objective}: understand ACL from scratch: view initial condition, add access
for one user, then create directory that inherits ACL automatically.

\textbf{Step 1}: Prepare file and view standard permissions without additional ACL.

\begin{lstlisting}[language=bash, caption={Preparing ACL test file}]
mkdir ~/lab-acl && cd ~/lab-acl
echo "Important data" > confidential.txt
chmod 640 confidential.txt
ls -l confidential.txt
getfacl confidential.txt
\end{lstlisting}

At this stage, \cmd{getfacl} only displays three basic entries: owner, group, and
others. No named user or named group yet.

\textbf{Step 2}: Give read access to one specific user without changing owner or group.

\begin{lstlisting}[language=bash, caption={Adding ACL for one user}]
setfacl -m u:userA:r confidential.txt
ls -l confidential.txt
getfacl confidential.txt
\end{lstlisting}

Notice two changes:
\begin{itemize}
    \item \cmd{ls -l} output shows \texttt{+} sign;
    \item \cmd{getfacl} output now has entry \texttt{user:userA:r--}.
\end{itemize}

\textbf{Step 3}: Create shared directory that inherits ACL to new files.

\begin{lstlisting}[language=bash, caption={Default ACL on directory}]
mkdir shared
setfacl -d -m u:userA:rwx shared
setfacl -d -m u:userB:r-x shared
getfacl shared

touch shared/inherited.txt
getfacl shared/inherited.txt
\end{lstlisting}

\textbf{Analysis}

\begin{enumerate}
    \item Why doesn't \cmd{getfacl confidential.txt} initially display specific users?
    \item After \cmd{setfacl -m u:userA:r confidential.txt}, what's the difference in
        \cmd{ls -l} and \cmd{getfacl} output?
    \item Why does file \texttt{inherited.txt} inherit ACL from directory \texttt{shared}?
\end{enumerate}

\textbf{Challenge}

Add one more ACL so that group \texttt{readonly-group} can only read
\file{confidential.txt}. After that, remove ACL for \texttt{userA} and verify the
final result with \cmd{getfacl}.

%% ========================================================================
\section{User and Group Management}
\label{sec:user-group-management}
%% ========================================================================

\textbf{Concept Explanation}

Linux stores user identity in three files: \file{/etc/passwd} (basic user data,
readable by everyone), \file{/etc/shadow} (password hash and aging policy, root only),
and \file{/etc/group} (group definitions). This system separates public information
from security secrets.

\textbf{Main Commands}

\begin{itemize}
    \item \texttt{useradd}, \texttt{usermod}, \texttt{userdel}: manage user accounts.
    \item \texttt{groupadd}, \texttt{groupmod}, \texttt{groupdel}: manage groups.
    \item \texttt{passwd}: change or lock password.
    \item \texttt{chage}: set password aging policy.
    \item \texttt{id}, \texttt{groups}: view user and group information.
\end{itemize}

\subsection{Reading User and Group Information}

\begin{lstlisting}[language=bash, caption={User and group query commands}]
id alice
groups alice
getent passwd alice
getent group developers
grep "^alice:" /etc/passwd
sudo grep "^alice:" /etc/shadow
\end{lstlisting}

\begin{notebox}
\cmd{id} is more accurate than \cmd{grep /etc/group} because it displays all user
groups, including primary group.
\end{notebox}

\subsection{User Lifecycle: Create, Modify, Delete}

\begin{lstlisting}[language=bash, caption={User account management}]
# create user with home directory and shell
sudo useradd -m -s /bin/bash alice
sudo passwd alice

# modify: add group, change shell, lock/unlock
sudo usermod -aG developers,sudo alice
sudo usermod -s /bin/zsh alice
sudo usermod -L alice
sudo usermod -U alice

# delete user (with or without home directory)
sudo userdel alice
sudo userdel -r alice
\end{lstlisting}

\textbf{Important options}:
\begin{itemize}
    \item \texttt{-m}: automatically create home directory
    \item \texttt{-s}: specify login shell
    \item \texttt{-aG}: add groups without removing old groups
    \item \texttt{-L}/\texttt{-U}: lock/unlock account
    \item \texttt{-r}: remove home directory when deleting user
\end{itemize}

\begin{warningbox}
Use \cmd{usermod -aG}, not \cmd{usermod -G}. Without \texttt{-a}, all old
supplementary groups will be replaced.
\end{warningbox}

\subsection*{Lab 9.3A --- Creating and Managing Users}

\textbf{Objective}: create new users, modify their properties, and understand the
difference between \cmd{useradd} and \cmd{usermod} options.

\begin{lstlisting}[language=bash, caption={User management practice}]
# create two users
sudo useradd -m -s /bin/bash userA
sudo useradd -m -s /bin/bash userB
sudo passwd userA
sudo passwd userB

# verify
id userA
getent passwd userA

# modify userA shell
sudo usermod -s /bin/zsh userA
getent passwd userA

# lock and unlock userB
sudo usermod -L userB
sudo passwd -S userB
sudo usermod -U userB
sudo passwd -S userB
\end{lstlisting}

\textbf{Questions}:
\begin{enumerate}
    \item What's the difference in \cmd{id userA} output before and after adding groups?
    \item How does \cmd{passwd -S userB} status change when account is locked?
\end{enumerate}

\subsection{Group Management}

\begin{lstlisting}[language=bash, caption={Group operations}]
sudo groupadd developers
sudo groupadd -g 2000 finance
sudo groupmod -n devteam developers
sudo gpasswd -a alice devteam
sudo gpasswd -d alice devteam
sudo groupdel devteam
\end{lstlisting}

\begin{notebox}
Group changes fully take effect after user logout-login or opening subshell with
\cmd{newgrp group-name}.
\end{notebox}

\subsection*{Lab 9.3B --- Group Management}

\textbf{Objective}: create groups, add users to groups, and verify membership.

\begin{lstlisting}[language=bash, caption={Group management practice}]
# create two groups
sudo groupadd labgroup
sudo groupadd readonly-group

# add userA to both groups
sudo usermod -aG labgroup,readonly-group userA

# add userB only to readonly-group
sudo usermod -aG readonly-group userB

# verify
id userA
id userB
getent group labgroup
getent group readonly-group
\end{lstlisting}

\textbf{Questions}:
\begin{enumerate}
    \item What does \cmd{id userA} display vs \cmd{groups userA}?
    \item Why is \texttt{-a} in \cmd{usermod -aG} important?
\end{enumerate}

\subsection{Password Policy and Aging}

\begin{lstlisting}[language=bash, caption={Password policy with chage}]
# view aging status
sudo chage -l alice

# set policy: password valid 90 days, warning 14 days before
sudo chage -M 90 -m 7 -W 14 alice

# lock account 30 days after password expires
sudo chage -I 30 alice

# set account expiration date
sudo chage -E 2025-12-31 alice

# force immediate password change
sudo chage -d 0 alice
sudo passwd -e alice
\end{lstlisting}

\begin{table}[htbp]
    \centering
    \caption{Most important \cmd{chage} options}
    \label{tab:chage-options}
    \begin{tblr}{
        colspec = {lcX[l]},
        width = \linewidth,
        hlines,
        vlines,
        row{1} = {font=\bfseries},
        cells = {font=\small}
    }
    \textbf{Option} & \textbf{Name} & \textbf{Meaning} \\
    \texttt{-l} & list & Display aging status \\
    \texttt{-M N} & maxdays & Password must be changed every N days \\
    \texttt{-m N} & mindays & Minimum N days before can change again \\
    \texttt{-W N} & warndays & Warning N days before expiry \\
    \texttt{-I N} & inactive & Lock account N days after expiry \\
    \texttt{-E DATE} & expire & Account expiration \\
    \end{tblr}
\end{table}

\subsection*{Lab 9.3C --- Password Aging Policy}

\textbf{Objective}: apply password aging policy and observe its effects.

\begin{lstlisting}[language=bash, caption={Password policy practice}]
# set aging policy for userA
sudo chage -M 60 -W 7 -m 1 userA
sudo chage -l userA

# force userA to change password on first login
sudo chage -d 0 userA

# lock userB password
sudo passwd -l userB
sudo passwd -S userB

# unlock back
sudo passwd -u userB
sudo passwd -S userB
\end{lstlisting}

\textbf{Questions}:
\begin{enumerate}
    \item What do the values displayed by \cmd{chage -l userA} mean?
    \item How can you prove userB is locked from \cmd{passwd -S} output?
    \item When should you use \cmd{chage -d 0} vs \cmd{passwd -e}?
\end{enumerate}

\textbf{Challenge}

Create a user named \texttt{intern} that:
\begin{itemize}
    \item has shell \texttt{/bin/bash};
    \item is member of \texttt{labgroup};
    \item is forced to change password on first login;
    \item has password expire after 45 days with 7 days warning before.
\end{itemize}

%% ========================================================================
\section{sudo and su Configuration}
\label{sec:sudo-su}
%% ========================================================================

\textbf{Concept Explanation}

\subsection{Difference Between \cmd{su} and \cmd{sudo}}

\cmd{su} switches user identity, while \cmd{sudo} runs specific commands with higher
privileges based on policy rules.

\begin{table}[htbp]
    \centering
    \caption{Comparison of \cmd{su} and \cmd{sudo}}
    \label{tab:su-vs-sudo}
    \begin{tblr}{
        colspec = {lX[l]X[l]},
        width = \linewidth,
        hlines,
        vlines,
        row{1} = {font=\bfseries},
        cells = {font=\small}
    }
    \textbf{Aspect} & \textbf{su} & \textbf{sudo} \\
    Authentication & Target user's password & Own password \\
    Scope & Switch fully to another session & Run specific commands \\
    Logging & Minimal / not default & Recorded in system log \\
    Granularity & Coarse & Very granular \\
    Modern practice & More limited & Recommended \\
    \end{tblr}
\end{table}

\begin{notebox}
In Ubuntu, the root account is usually disabled for direct login; therefore \cmd{sudo}
becomes the standard administration mechanism.
\end{notebox}

\textbf{Commands Used}

\begin{itemize}
    \item \texttt{su - user}: switch to another user with full login environment.
    \item \texttt{sudo command}: run one command with higher privileges.
    \item \texttt{sudo -i}: open root shell with login environment.
    \item \texttt{sudo -l}: view available sudo rights.
    \item \texttt{visudo}: edit sudoers rules with syntax validation.
\end{itemize}

\subsection{Basic Usage of \cmd{su} and \cmd{sudo}}

\begin{lstlisting}[language=bash, caption={Most commonly used su and sudo commands}]
# su
su - alice
su - alice -c "ls -la /home/alice"

# sudo
sudo apt update
sudo -u alice cat /home/alice/report.txt
sudo -i
sudo -s
sudo -l
sudo -k
sudo -e /etc/hosts
\end{lstlisting}

Brief guidelines:

\begin{itemize}
    \item use \cmd{su - user}, not \cmd{su user}, if you must switch sessions;
    \item use \cmd{sudo -i} for root shell with full login environment;
    \item use \cmd{sudo -e} or \cmd{sudoedit} to edit system files more securely.
\end{itemize}

\subsection{Safe sudoers Configuration}

Main sudo configuration is in \file{/etc/sudoers}, but best practice is to store
additional rules in \dir{/etc/sudoers.d/}.

\begin{warningbox}
Always edit sudo rules with \cmd{visudo}. Syntax errors can lock you out of
administrative access to the system.
\end{warningbox}

Basic rule format:

\begin{center}
\texttt{WHO\ \ WHERE=(AS\_WHO)\ \ COMMAND}
\end{center}

\begin{lstlisting}[language=bash, caption={Representative sudoers rule examples}]
# full access for specific user
alice ALL=(ALL:ALL) ALL

# apt access without password
bob ALL=(root) NOPASSWD: /usr/bin/apt

# default Ubuntu sudo group
%sudo ALL=(ALL:ALL) ALL

# limited permission for nginx service
%webadmin ALL=(root) /bin/systemctl start nginx, \
                     /bin/systemctl stop nginx, \
                     /bin/systemctl restart nginx, \
                     /bin/systemctl status nginx
\end{lstlisting}

\begin{lstlisting}[language=bash, caption={Managing sudoers files}]
sudo visudo
sudo visudo -f /etc/sudoers.d/developers
sudo visudo -c
sudo grep sudo /var/log/auth.log | tail -20
\end{lstlisting}

\subsection*{Lab 9.4 --- sudo Configuration}

\textbf{Objective}: create limited sudo rules, verify access rights, and read logs.

\textbf{Step 1}: Create special sudo configuration file for \texttt{userA}.

\begin{lstlisting}[language=bash, caption={Creating separate sudoers file}]
sudo visudo -f /etc/sudoers.d/lab-userA
\end{lstlisting}

This command opens a safe editor specifically for new sudoers file. If syntax is wrong,
\cmd{visudo} will warn before saving.

Fill file with following rules:

\begin{lstlisting}[language=bash, caption={sudo configuration content for userA}]
userA ALL=(root) NOPASSWD: /usr/bin/apt update, /usr/bin/apt upgrade
userA ALL=(root) /bin/systemctl status *
\end{lstlisting}

First line means \texttt{userA} can run two \cmd{apt} commands without password.
Second line means \texttt{userA} can view any service status, but still follows
normal authentication policy.

\textbf{Step 2}: Verify active rules and test results.

\begin{lstlisting}[language=bash, caption={Verify sudo rights and log}]
sudo -l -U userA
sudo grep "userA" /var/log/auth.log | tail -10
\end{lstlisting}

\cmd{sudo -l -U userA} is used to check active rules from \texttt{userA} account
perspective. Log in \file{/var/log/auth.log} helps verify that sudo usage is actually
recorded.

\textbf{Analysis}

\begin{enumerate}
    \item Why are rules stored in \dir{/etc/sudoers.d/}, not directly in \file{/etc/sudoers}?
    \item Which commands can be run without password, and which still need authentication?
    \item What information is recorded in sudo log?
\end{enumerate}

\textbf{Challenge}

Add one new rule so \texttt{userA} can run \cmd{/bin/systemctl restart ssh} but
cannot run \cmd{reboot}.

%% ========================================================================
\section{Disk Quota}
\label{sec:disk-quota}
%% ========================================================================

\textbf{Concept Explanation}

Quota limits disk space usage and/or number of files owned by user or group. Its
function is simple: prevent one user from consuming all storage resources.

\subsection{Disk Quota Concept}
\label{sec:disk-quota-konsep}

There are two quota dimensions:

\begin{description}
    \item[Block quota] limits total disk space used. In Linux, quota tools generally
        use kilobyte units.
    \item[Inode quota] limits number of files/directories, regardless of their size.
\end{description}

Additionally, there are three policy terms that must be understood:

\begin{table}[htbp]
    \centering
    \caption{Types of limits in disk quota}
    \label{tab:quota-limits}
    \begin{tblr}{
        colspec = {lQ[2cm,l]X[l]},
        width = \linewidth,
        hlines,
        vlines,
        row{1} = {font=\bfseries},
        cells = {font=\small}
    }
    \textbf{Type} & \textbf{Applies to} & \textbf{Meaning} \\
    Soft limit & Blocks and inodes & Can be temporarily exceeded during grace period \\
    Hard limit & Blocks and inodes & Cannot be exceeded at all \\
    Grace period & Soft limit & Tolerance time before soft limit is treated like hard limit \\
    \end{tblr}
\end{table}

\begin{notebox}
Inode quota is important when users create very many small files. Disk space may still
be available, but filesystem inodes can run out first.
\end{notebox}

\textbf{Commands Used}

\begin{itemize}
    \item \texttt{quotacheck}: create or update quota database.
    \item \texttt{quotaon}/\texttt{quotaoff}: enable or disable enforcement.
    \item \texttt{edquota}: edit quota limits interactively.
    \item \texttt{setquota}: set quota limits from command line.
    \item \texttt{quota}, \texttt{repquota}: view quota usage.
\end{itemize}

\subsection{Enabling Quota}

Quota activation workflow usually consists of four steps:

\begin{enumerate}
    \item install quota package,
    \item enable \texttt{usrquota} and/or \texttt{grpquota} options on filesystem,
    \item create quota database with \cmd{quotacheck},
    \item enable enforcement with \cmd{quotaon}.
\end{enumerate}

\begin{lstlisting}[language=bash, caption={Basic quota activation workflow}]
sudo apt install -y quota quotatool

# /etc/fstab
# UUID=abc123 /home ext4 defaults,usrquota,grpquota 0 2

sudo mount -o remount /home
sudo quotacheck -cugm /home
sudo quotaon -v /home
sudo repquota -s /home
\end{lstlisting}

\subsection{Setting and Viewing Quota}

The most important commands are:

\begin{lstlisting}[language=bash, caption={Core quota management commands}]
# interactive editor
sudo edquota -u alice
sudo edquota -g developers

# without editor
sudo setquota -u alice 1048576 1572864 10000 15000 /home

# view status
quota -s
sudo quota -u alice
sudo repquota -s /home

# grace period
sudo edquota -t
\end{lstlisting}

Meaning of arguments in \cmd{setquota -u alice 1048576 1572864 10000 15000 /home}:

\begin{itemize}
    \item soft block = 1~GB
    \item hard block = 1.5~GB
    \item soft inode = 10,000
    \item hard inode = 15,000
\end{itemize}

\subsection{Advanced Operations to Know}

\begin{lstlisting}[language=bash, caption={Advanced quota operations}]
# copy quota from one user to another
sudo edquota -p alice bob charlie

# send warning email
sudo warnquota

# turn off quota
sudo quotaoff -v /home
\end{lstlisting}

\subsection*{Lab 9.5 --- Disk Quota}

\textbf{Objective}: understand quota flow safely on loopback filesystem without
modifying main filesystem.

\textbf{Step 1}: Create small filesystem image and mount with quota options.

\begin{lstlisting}[language=bash, caption={Preparing quota test filesystem}]
sudo dd if=/dev/zero of=/tmp/quota-test.img bs=1M count=100
sudo mkfs.ext4 /tmp/quota-test.img
sudo mkdir -p /mnt/quota-test
sudo mount -o loop,usrquota,grpquota /tmp/quota-test.img /mnt/quota-test
\end{lstlisting}

Image file is used to make the lab safe: you don't need to modify main filesystem
like \dir{/home}. Options \texttt{usrquota,grpquota} enable both quota types simultaneously.

\textbf{Step 2}: Create quota database and enable enforcement.

\begin{lstlisting}[language=bash, caption={Initialization and quota activation}]
sudo quotacheck -cug /mnt/quota-test
sudo quotaon -v /mnt/quota-test
sudo repquota /mnt/quota-test
\end{lstlisting}

\cmd{quotacheck -cug} creates user and group quota database. After that, \cmd{quotaon}
enables enforcement, and \cmd{repquota} displays initial report.

\textbf{Step 3}: Set quota for test user and observe results.

\begin{lstlisting}[language=bash, caption={Setting quota for userA}]
sudo edquota -u userA
# example: soft block 5120, hard block 10240
sudo repquota /mnt/quota-test
\end{lstlisting}

Values above use KB units. So \texttt{5120} means about 5 MB, and \texttt{10240}
means about 10 MB.

\textbf{Step 4}: Clean up test environment after completion.

\begin{lstlisting}[language=bash, caption={Quota test cleanup}]
sudo quotaoff /mnt/quota-test
sudo umount /mnt/quota-test
sudo rm /tmp/quota-test.img
\end{lstlisting}

\textbf{Analysis}

\begin{enumerate}
    \item What's the difference between soft limit and hard limit when quota starts being exceeded?
    \item Why does this lab use loopback filesystem, not directly \dir{/home}?
    \item From \cmd{repquota} output, what information shows quota is already active?
\end{enumerate}

\textbf{Challenge}

Try setting new quota for \texttt{userA} with very small inode limit, then explain
when inode restriction is more important than block restriction.

%% ========================================================================
\section{Summary}
\label{sec:file-user-summary}
%% ========================================================================

This chapter discussed five core concept groups:

\begin{description}
    \item[Permissions] owner, group, others; \texttt{rwx} permissions; \cmd{chmod},
        \cmd{chown}, \cmd{chgrp}; special bits; and umask.
    \item[ACL] extension of Unix permissions for granular per-user or per-group
        access needs, including directory default ACL and mask.
    \item[User/Group] user identity stored in \file{/etc/passwd}, \file{/etc/shadow},
        and \file{/etc/group}; managed with command families \cmd{useradd},
        \cmd{usermod}, \cmd{userdel}, \cmd{groupadd}, and \cmd{gpasswd}.
    \item[sudo and su] \cmd{sudo} is a delegation mechanism for administrative rights
        that is safer and more auditable than \cmd{su}.
    \item[Disk Quota] block and inode restrictions through soft limit, hard limit,
        grace period, and commands like \cmd{edquota}, \cmd{setquota}, and \cmd{repquota}.
\end{description}

%% ========================================================================
\section{Exercises}
\label{sec:file-user-exercises}
%% ========================================================================

\begin{exercisebox}[Exercise 9.A --- Audit and Collaboration]
\begin{enumerate}
    \item Find active SUID files with \cmd{find / -perm -4000 -type f 2>/dev/null},
        then explain three files you recognize with their reasons.
    \item Find \textit{world-writable} directories and determine which are valid and
        which are risky.
    \item Design standard permission configuration and ACL for project directory
        \dir{/srv/webapp} so group \texttt{webapp-team} can write, user
        \texttt{deploy} can only read, and new files always inherit project group.
\end{enumerate}
\end{exercisebox}

\begin{exercisebox}[Exercise 9.B --- Account Policy and Quota]
Write steps to create user \texttt{intern}, add to group \texttt{labgroup},
force password change every 45 days (warning 7 days), give sudo permission only
for \cmd{systemctl status}, and set simple space and inode quota on \dir{/home}.
\end{exercisebox}

%% ========================================================================
